% O014 / D80 — Methods of Algebra, Volume 2 — id-ID
% Sumber: appendix2.tex, baris 1--137 (inklusif)
% ID stabil unit: o014.aljabr2.appendix-b.profinite-groups
% Sumber asli © 2024 Wen-Wei Li, CC BY 4.0.

% Untuk disertakan

\chapter{Pengantar objek-ind dan objek-pro}\label{sec:app-Ind}
Untuk suatu kategori $\mathcal{C}$, apakah yang dimaksud dengan objek-ind di atasnya? Melalui pembenaman Yoneda $\mathcal{C} \to \mathcal{C}^{\wedge}$, objek-ind dapat didefinisikan secara ketat sebagai objek dalam $\mathcal{C}^\wedge$ yang berbentuk $X = \Yinjlim X_i$, dengan indeks $i$ menjelajahi suatu kategori kecil terfilter $I$ dan $X_i \in \Obj(\mathcal{C})$; simbol $\Yinjlim$ dipakai untuk membedakan $\varinjlim$ dalam $\mathcal{C}^{\wedge}$ dari yang ada dalam $\mathcal{C}$. Dapat dibuktikan bahwa
\[ \Hom(X, Y) = \varprojlim_i \varinjlim_j \Hom_{\mathcal{C}}(X_i, Y_j), \quad X = \Yinjlim X_i, \; Y = \Yinjlim Y_j. \]
Karena itu, secara kasar objek-ind $X$ dapat dibayangkan sebagai suatu ``sistem terfilter'', yang direpresentasikan oleh keluarga objek $X_i$ dalam $\mathcal{C}$ bersama keluarga morfisme kompatibel $X_i \to X_j$, dengan $i \to j$ menjelajahi $\Mor(I)$; tetapi pilihan $(X_i)_i$ yang merepresentasikan $X$ tidak tunggal. Konsep ini juga mempunyai versi dual yang disebut objek-pro di atas $\mathcal{C}$. Penyajian di sini mengikuti pendekatan \cite{KS06}.

Sebagai sumber contoh, \S\ref{sec:profinite-groups} dimulai dengan grup profinit: grup-grup ini sekaligus merupakan kelas khusus grup topologis dan dapat dicirikan secara ekuivalen sebagai $\varprojlim$ terfilter dari grup hingga dalam kategori grup topologis (Catatan \ref{rem:profinite-group-lim}). Berdasarkan ini, Contoh \ref{eg:profinite-group} kelak menunjukkan bahwa grup profinit tidak lain adalah objek-pro di atas kategori grup hingga, sehingga dapat diperlakukan dengan metode aljabar. Kelas grup topologis ini sangat sering digunakan dalam matematika dan juga menjadi pengetahuan latar bagi beberapa bab buku ini.

Bagian~\ref{sec:ind-pro} memberikan definisi formal objek-ind dan objek-pro,
deskripsi morfismenya, dan beberapa contoh awal. Pembahasan selanjutnya
berfokus pada objek-ind; \S\ref{sec:Indization} mengkaji kategori
$\IndC\mathcal{C}$ yang terdiri atas semua objek-ind di atas $\mathcal{C}$,
disebut pula ind-isasi dari $\mathcal{C}$, dan memuat $\mathcal{C}$ sebagai
subkategori penuh. Kita akan mengkaji cara mengenali objek-ind dalam
$\mathcal{C}^{\wedge}$ (Proposisi \ref{prop:recognize-ind}), yakni mencirikan
keterwakilan-ind suatu funktor $\mathcal{C}^{\opp} \to \cate{Set}$, lalu
membahas keterkaitan antara $\mathcal{C} \to \IndC\mathcal{C}$ dan limit;
khususnya, kita akan menunjukkan bahwa $\IndC\mathcal{C}$ mempunyai semua
$\varinjlim$ kecil terfilter.

Secara dual, $\mathcal{C}$ juga mempunyai pro-isasi: semua objek-pro membentuk kategori $\ProC\mathcal{C} = \IndC\left(\mathcal{C}^{\opp}\right)^{\opp}$.

Ind-isasi funktor merupakan pokok bahasan \S\ref{sec:Indization-functor}. Hasilnya mencakup cara memperluas funktor $\mathcal{C} \to \mathcal{D}$ ke $\IndC\mathcal{C}$ (Definisi--Proposisi \ref{def:Ind-hat-functor}) serta cara mengenali apakah suatu kategori $\mathcal{C}$ merupakan ind-isasi dari subkategori penuh $\mathcal{C}'$ (Proposisi \ref{prop:Ind-hat-small}).

Dalam \S\ref{sec:Indization-Abel}, hasil-hasil ini diterapkan pada kategori abelian $\mathcal{C}$ untuk menunjukkan bahwa kedua pembenaman $\mathcal{C} \to \IndC\mathcal{C}$ dan $\mathcal{C} \to \ProC\mathcal{C}$ merupakan funktor eksak setia-penuh di antara kategori-kategori abelian (Teorema \ref{prop:Ind-Abel}); di sana kita juga membahas keeksakan perluasan funktor pada kategori abelian (Proposisi \ref{prop:Ind-ext-exactness}). Lebih lanjut, Teorema \ref{prop:Ind-Grothendieck} menunjukkan bahwa ind-isasi suatu kategori abelian kecil adalah kategori Grothendieck.

Berdasarkan hasil-hasil tersebut dan fakta dasar bahwa kategori Grothendieck mempunyai kopembangun injektif, \S\ref{sec:FM} akan membuktikan teorema pembenaman Freyd--Mitchell: setiap kategori abelian kecil dapat dibenamkan secara eksak dan setia-penuh ke dalam $\cated{Mod}R$, dengan $R$ suatu gelanggang yang direalisasikan pada himpunan kecil. Pendekatan melalui ind-isasi bukanlah yang paling elementer ataupun yang paling singkat; namun, dibandingkan dengan teorema Freyd--Mitchell itu sendiri, ind-isasi boleh jadi merupakan perangkat yang lebih berguna.

\section{Pendahuluan: grup profinit}\label{sec:profinite-groups}
Bagian ini mula-mula mendefinisikan grup profinit dari sudut pandang topologis: grup-grup ini adalah grup topologis dengan sifat tertentu. Dalam sedikit keadaan ketika ukuran himpunan harus dibedakan, secara bawaan grup topologis yang kita tinjau direalisasikan pada himpunan kecil; selain itu hal ini tidak akan disebutkan lagi.

\begin{definition}\label{def:profinite-group}
	\index{grup profinit (profinite group)}
	Sebuah \emph{grup profinit} adalah grup topologis $G$ yang memenuhi syarat-syarat berikut:
	\begin{itemize}
		\item $G$ adalah grup Hausdorff kompak;
		\item $G$ mempunyai basis lingkungan di unsur identitas $1_G$ yang terdiri atas subgrup normal.
	\end{itemize}
	
	Nyatakan kategori semua grup topologis dengan $\cate{TopGrp}$, dengan morfisme berupa homomorfisme kontinu. Semua grup profinit lalu membentuk subkategori penuh dari $\cate{TopGrp}$.
\end{definition}

Jika $G$ adalah grup profinit, basis lingkungan di $1_G$ dapat dipilih sebagai semua subgrup normal buka $K \lhd G$. Untuk pembahasan terperinci lihat \cite[\S 1.9]{FL14} atau \cite[\S 4.10]{Li1}.

\begin{remark}
	Pencirian lain berbunyi: grup topologis $G$ adalah profinit jika dan hanya jika ia merupakan grup kompak yang tak terhubung total; lihat \cite[\S 1.7 dan Proposisi 1.9.3]{FL14}. Fakta terkait berikut juga berguna: ruang topologis Hausdorff $E$ bersifat kompak lokal dan tak terhubung total jika dan hanya jika setiap $e \in E$ mempunyai basis lingkungan yang terdiri atas subhimpunan kompak buka.
\end{remark}

Jika $H$ adalah subgrup buka dari $G$, pilih wakil $g$ untuk setiap koset
$\bar{g} \in G/H$. Dari
$G = \bigsqcup_{\bar{g} \in G/H} gH$ dan kekompakan $G$, diperoleh bahwa
$(G:H)$ hingga. Selain itu,
$H = G \smallsetminus \bigcup_{\bar{g} \neq H} gH$ menunjukkan bahwa $H$
tertutup.

Kita catat beberapa sifat dasar lebih lanjut.

\begin{lemma}
	Misalkan $G$ grup topologis dan $H$ subgrupnya. Maka:
	\begin{enumerate}[(i)]
		\item pemetaan hasil bagi $\pi: G \to G/H$ bersifat terbuka terhadap topologi hasil bagi pada $G/H$;
		\item $G/H$ dengan topologi hasil bagi merupakan ruang Hausdorff jika dan hanya jika $H$ tertutup;
		\item topologi hasil bagi pada $G/H$ bersifat diskret jika dan hanya jika $H$ terbuka.
	\end{enumerate}
	Tentu saja pernyataan-pernyataan di atas juga berlaku bagi $H \backslash G$.
\end{lemma}
\begin{proof}
	Untuk (i), jika $U$ adalah subhimpunan terbuka dari $G$, maka $\pi^{-1}(\pi(U)) = \bigcup_{h \in H} Uh$ juga terbuka; jadi, menurut definisi topologi hasil bagi, $\pi(U)$ terbuka.
	
	Arah ``hanya jika'' pada (ii) mengikuti dari $H = \pi^{-1}(1_G \cdot H)$. Untuk arah ``jika'', diberikan koset berbeda $xH \neq yH$, pilih subhimpunan terbuka $V \ni 1_G$ dari $G$ sedemikian sehingga $Vx \cap yH = \emptyset$, lalu pilih subhimpunan terbuka $U \ni 1_G$ sedemikian sehingga $U^{-1} U \subset V$. Ini menjamin $UxH \cap UyH = \emptyset$. Sementara itu, (i) menunjukkan bahwa $\pi(Ux)$ dan $\pi(Uy)$ sama-sama terbuka.
	
	Arah ``hanya jika'' pada (iii) kembali mengikuti dari $H = \pi^{-1}(1_G \cdot H)$. Sekarang andaikan $H$ terbuka. Maka (i) menyiratkan bahwa $\{1_G \cdot H\} = \pi(H)$ terbuka; melalui translasi, semua himpunan singleton dalam $G/H$ terbuka. Ini membuktikan arah ``jika''.
\end{proof}

\begin{proposition}
	Misalkan $G$ grup profinit.
	\begin{enumerate}[(i)]
		\item Jika $H$ adalah subgrup tertutup dari $G$, maka $H$ merupakan grup profinit.
		\item Jika $H$ adalah subgrup normal tertutup dari $G$, maka $G/H$ dengan topologi hasil bagi merupakan grup profinit.
	\end{enumerate}
\end{proposition}
\begin{proof}
	Untuk (i), subgrup tertutup $H$ tentu kompak dan Hausdorff; suatu basis lingkungan di $1_G$ diberikan oleh semua $K \cap H$, dengan $K$ menjelajahi subgrup normal buka dari $G$.
	
	Untuk (ii), kita tahu bahwa ketertutupan $H$ menyiratkan $G/H$ Hausdorff, sedangkan kekompakan $G$ menyiratkan $G/H$ kompak. Basis lingkungan bagi $G/H$ di $1_G \cdot H$ diberikan oleh semua $KH/H$, dengan $K$ menjelajahi subgrup normal buka dari $G$; kita juga tahu bahwa $KH/H$ merupakan subgrup normal buka dari $G/H$.
\end{proof}

Mudah dibuktikan bahwa produk suatu keluarga grup profinit $(G_i)_i$, dengan $i$ menjelajahi suatu himpunan kecil, tetap profinit. Untuk setiap kategori kecil $I$ dan funktor $\beta: I^{\opp} \to \cate{TopGrp}$, yang secara singkat ditulis $(G_i)_{i \in \Obj(I)}$ (dengan $G_i := \beta(i)$), kita membekali
\[ \left( \varprojlim_i G_i := \varprojlim \beta \right) \hookrightarrow \prod_i G_i \]
dengan topologi yang diwarisi ruas kanan dari ruang produk, sehingga ia menjadi subgrup tertutup. Ini memberikan $\varprojlim$ dalam $\cate{TopGrp}$. Jadi, jika setiap $G_i$ profinit, maka $\varprojlim_i G_i$ juga profinit. Sebagai akibatnya, kategori semua grup profinit merupakan kategori lengkap.

\begin{example}
	Contoh paling sederhana adalah ketika $G$ merupakan grup hingga; dalam hal ini hanya ada satu pilihan topologi Hausdorff, yakni topologi diskret. Contoh-contoh berikut juga lazim.
	\begin{itemize}
		\item Ambil grup Galois $G := \Gal(E|F)$ dari ekstensi Galois $E|F$, dengan topologi Krull \cite[\S 9.2]{Li1}. Basis lingkungan yang bersesuaian terdiri atas $\Gal(E|L)$, dengan $L|F$ menjelajahi subekstensi Galois hingga dari $E|F$.
		
		\item Untuk suatu bilangan prima $p$, nyatakan gelanggang bilangan bulat $p$-adik dengan $\Z_p$, dan nyatakan grup matriks $n \times n$ berentri di $\Z_p$ dengan determinan $1$ sebagai $\SL(n, \Z_p)$. Dengan topologi yang diwarisi dari ruang matriks $n \times n$, $\mathrm{M}_n(\Z_p)$, grup ini juga profinit.
		
		\item Untuk sebarang grup $G$, pelengkapan profinitnya didefinisikan sebagai
		\[ \hat{G} := \varprojlim_{\substack{N \lhd G \\ (G:N) < \infty}} G/N, \]
		dengan setiap $G/N$ diberi topologi diskret. Sesuai namanya, ini merupakan grup profinit.
	\end{itemize}
\end{example}

\begin{remark}\label{rem:profinite-group-lim}
	Pembahasan sebelumnya menunjukkan bahwa apabila grup hingga diberi topologi diskret, maka $\varprojlim$ kecilnya merupakan grup profinit. Sebaliknya, setiap grup profinit $G$ dapat ditulis sebagai $\varprojlim_i G_i$, dengan $(I, \leq)$ suatu himpunan kecil terurut parsial terfilter dan setiap $G_i$ grup hingga; lihat \cite[Teorema 4.10.6]{Li1}. Secara konkret, kita dapat mengambil basis lingkungan $(K_i)_{i \in I}$ di $1_G$ yang terdiri atas subgrup normal buka, mengurutkan himpunan indeks secara terbalik terhadap inklusi, dan menetapkan $G_i := G/K_i$.
	
	Sebagai contoh grup Galois, $\Gal(E|F) = \varprojlim_{L|F} \Gal(L|F)$, dengan $L|F$ menjelajahi subekstensi Galois hingga dari $E|F$.
	
	Tinjau grup profinit $G = \varprojlim_i G_i$ dan $H = \varprojlim_j H_j$ dengan penyajian seperti di atas. Karena semua $\Ker[G \to G_i]$ membentuk basis lingkungan di $1_G$, kita memperoleh isomorfisme kanonik
	\begin{align*}
		\Hom_{\cate{TopGrp}}(G, H) & = \Hom_{\cate{TopGrp}}\left(\varprojlim_i G_i, \varprojlim_j H_j \right) \\
		& \simeq \varprojlim_j \Hom_{\cate{TopGrp}}\left( \varprojlim_i G_i, H_j \right) \quad \text{($\because\; $ sifat universal $\varprojlim$)} \\
		& \simeq \varprojlim_j \varinjlim_i \Hom_{\cate{Grp}}(G_i, H_j) \quad \text{($\because\; $ setiap homomorfisme kontinu memfaktor melalui suatu $G_i$)};
	\end{align*}
	Atas dasar ini, grup profinit juga dapat diperlakukan secara aljabar; lihat Contoh \ref{eg:profinite-group} di bawah.
\end{remark}

\begin{definition}
	\index{grup pro-$p$}
	Misalkan $p$ bilangan prima. Sebuah grup profinit $G$ disebut \emph{grup pro-$p$} jika mempunyai sifat berikut: untuk setiap subgrup normal buka $K \lhd G$ yang cukup kecil, grup hasil bagi $G/K$ merupakan grup-$p$. Secara ekuivalen, $G$ dapat ditulis sebagai $\varprojlim_i G_i$, dengan $(I, \leq)$ poset terfilter dan semua $G_i$ grup-$p$.
\end{definition}

Kelas grup ini sering muncul dalam teori bilangan. Latihan memuat pembahasan lebih lanjut mengenai grup pro-$p$.

Hasil berikut menunjukkan bahwa pemetaan hasil bagi grup profinit oleh subgrup tertutup mempunyai perilaku topologis yang sederhana, sangat berbeda dari keadaan grup Lie yang lazim.

\begin{lemma}\label{prop:profinite-section}
	Misalkan $H$ dan $H'$ subgrup tertutup dari grup profinit $G$, dengan $H' \subset H$. Maka pemetaan hasil bagi $\pi: G/H' \twoheadrightarrow G/H$ mempunyai penampang kontinu $s: G/H \to G/H'$; dengan kata lain, $s$ adalah pemetaan kontinu yang memenuhi $\pi s = \identity_{G/H}$.
	
	Pernyataan serupa berlaku bagi pemetaan hasil bagi $H' \backslash G \twoheadrightarrow H \backslash G$.
\end{lemma}
\begin{proof}
	Mula-mula tangani kasus ketika $(H:H')$ hingga. Dalam hal ini $H$ merupakan gabungan saling lepas dari $H'$ dan sejumlah hingga translasi kirinya, sehingga $H'$ terbuka dalam $H$. Terdapat subgrup normal buka $U \lhd G$ sedemikian sehingga $U \cap H \subset H'$. Maka pembatasan $\pi$ merupakan bijeksi $UH'/H' \to UH/H$; oleh kekompakan, ini merupakan homeomorfisme, dan inversnya memberikan penampang kontinu pada subhimpunan terbuka $UH/H$ dari $G/H$. Akhirnya, dengan translasi kita perluas penampang itu ke seluruh $G/H$.
	
	Untuk kasus umum, pembaca diminta mereduksi masalah ke kasus $H' = \{1_G\}$. Dalam kasus ini, tinjau semua data $(T,t)$, dengan $T \subset H$ subgrup tertutup dan $t$ penampang kontinu dari $G/T \twoheadrightarrow G/H$. Beri data-data tersebut urutan parsial berikut: $(T_1,t_1) \preceq (T_2,t_2)$ berarti $T_2 \subset T_1$ dan $t_1$ memfaktor sebagai $G/H \xrightarrow{t_2} G/T_2 \twoheadrightarrow G/T_1$. Lema Zorn menjamin adanya unsur maksimal. Pokoknya, jika $((S_i,s_i))_i$ adalah rantai terhadap $\preceq$ dan $S' := \bigcap_i S_i$, maka pemetaan kanonik $G/S' \to \varprojlim_i G/S_i$ kontinu, injektif, dan bercitra rapat; kekompakan lalu menyiratkan bahwa pemetaan itu adalah homeomorfisme. Karena itu dapat didefinisikan $s' := \varprojlim_i s_i: G/H \to G/S'$ sehingga $(S',s')$ menjadi batas atas rantai tersebut.
	
	Tinjau suatu unsur maksimal $(S,s)$ dari poset di atas. Jika $S=\{1\}$, pembuktian selesai. Jika $S \neq \{1\}$, pilih subgrup normal buka $U$ dari $G$ sedemikian sehingga $S \cap U \neq S$. Karena $(S:S\cap U)$ hingga, langkah pertama pembuktian memberikan penampang kontinu dari $G/(S\cap U) \to G/S$; setelah dikomposisikan dengan $s$, diperoleh penampang kontinu dari $G/(S\cap U) \to G/H$. Ini bertentangan dengan kemaksimalan $(S,s)$.
	
	Terakhir, versi untuk $H' \backslash G \twoheadrightarrow H \backslash G$ sama persis.
\end{proof}

Kasus khusus terpenting adalah $H'=\{1_G\}$. Dalam hal ini penampang kontinu $s$ menginduksi homeomorfisme ruang topologis $(G/H) \times H \rightiso G$ yang memetakan $(x,h)$ ke $s(x)h$. Untuk $G \twoheadrightarrow H \backslash G$, kesimpulannya tentu serupa.
